\documentclass[11pt]{article}
\usepackage[a4paper,margin=2.4cm]{geometry}
\usepackage[english]{babel}
\usepackage[utf8]{inputenc}
\usepackage[T1]{fontenc}
\usepackage{lmodern}
\usepackage{microtype}
\usepackage{graphicx}
\usepackage{xcolor}
\usepackage{tcolorbox}
\usepackage{booktabs}
\usepackage{enumitem}
\usepackage{hyperref}
\usepackage{amsmath,amssymb}
\usepackage{titlesec}

\definecolor{accent}{HTML}{4477AA}
\definecolor{boxbg}{HTML}{F5F5F0}
\definecolor{boxrule}{HTML}{888888}
\definecolor{questbg}{HTML}{FFF7E6}
\definecolor{questrule}{HTML}{CC8800}
\definecolor{warnbg}{HTML}{FCEEE6}
\definecolor{warnrule}{HTML}{B8553A}

\hypersetup{
  colorlinks=true,
  linkcolor=accent,
  urlcolor=accent,
  citecolor=accent,
}

\graphicspath{{../figures/}}

\titleformat{\section}{\Large\bfseries\color{accent}}{\thesection.}{0.6em}{}
\titleformat{\subsection}{\large\bfseries}{\thesubsection.}{0.5em}{}

\newtcolorbox{cmd}{
  colback=boxbg, colframe=boxrule, boxrule=0.4pt, arc=2pt,
  fontupper=\ttfamily\small, left=8pt, right=8pt, top=4pt, bottom=4pt,
  before skip=8pt, after skip=8pt
}
\newtcolorbox{quest}[1][Question]{
  colback=questbg, colframe=questrule, boxrule=0.5pt, arc=2pt,
  title=#1, fonttitle=\bfseries\sffamily,
  left=8pt, right=8pt, top=4pt, bottom=4pt,
  before skip=8pt, after skip=8pt
}
\newtcolorbox{warn}[1][Note]{
  colback=warnbg, colframe=warnrule, boxrule=0.5pt, arc=2pt,
  title=#1, fonttitle=\bfseries\sffamily,
  left=8pt, right=8pt, top=4pt, bottom=4pt,
  before skip=8pt, after skip=8pt
}

\title{\textbf{\texttt{impop\_k} hands-on tutorial}\\[0.4em]
       \large IBS, IBD, and local ancestry from pangenome alignments}
\author{MEMPANG workshop}
\date{}

\begin{document}
\maketitle

\noindent
\textbf{Time budget.} 45 minutes. Compute time across the entire tutorial is
a few seconds; the rest of the time is reading outputs, plotting, and
answering interpretation questions. Work alone or in pairs.

\vspace{0.4em}
\noindent
\textbf{Setup assumptions.} You have the \texttt{workshop/} directory on disk
(cloned from the course repository), the four pre-built binaries in
\texttt{workshop/bin/}, the precomputed identity matrices in
\texttt{workshop/data/}, and Python 3 with \texttt{numpy} and
\texttt{matplotlib}. All commands below \emph{must} be run from the
\texttt{workshop/} directory.

\vspace{0.4em}
\noindent
\textbf{Where things came from.} Yesterday you ran \texttt{impg similarity}
on PAF alignments to obtain pairwise identity per window. The two TSVs in
\texttt{data/} are exactly that --- precomputed on \texttt{chr12}, so you
don't pay the cost again today. The reproduction script is
\texttt{code/regenerate\_data.sh}.

\vspace{0.4em}
\noindent
\textbf{Why chr12 long arm.} Every section below works on the long (q) arm
of \texttt{chr12} --- roughly 40 to 130 Mb --- staying clear of the
centromere at $\sim$34--38 Mb and of both telomeres. The arm gives us
enough length to resolve recombination breakpoints visually.

\section{Setup (3 min)}

\begin{cmd}
\$ cd workshop \\
\$ ls bin/ data/ code/ \\
\$ ./bin/ancestry --help | head -3
\end{cmd}

You should see four binaries (\texttt{ibs}, \texttt{ibd}, \texttt{ancestry},
\texttt{jacquard}), TSV/text files in \texttt{data/}, and Python/Bash scripts
in \texttt{code/}. The third command should print three lines of help text
starting with ``Local ancestry inference''.

\begin{warn}[If something fails]
If \texttt{./bin/ancestry --help} reports a missing library, your VM's
\texttt{glibc} is older than the one the binaries were built against
(2.39). Tell the instructor; the fix is to copy the binaries from a Linux
laptop with a more recent \texttt{glibc} or to build them from source
(\texttt{cargo build --release} in the impop repository).
\end{warn}

\section{Part 1 --- IBS: looking at identity windows (7 min)}

\subsection*{The dataset}

\texttt{data/identity\_chr12\_pedigree.tsv} contains pairwise identity per
10-kb window for the CEPH 1463 platinum pedigree on
\texttt{chr12:40--130\,Mb} (the long q arm). Eighteen haplotypes total:
five full siblings (NA12879, NA12881, NA12882, NA12885, NA12886), the
two paternal grandparents (NA12889, NA12890), and the two maternal
grandparents (NA12891, NA12892). The two grandparental couples are
mutually unrelated.

\begin{cmd}
\$ head -3 data/identity\_chr12\_pedigree.tsv \\
\$ wc -l data/identity\_chr12\_pedigree.tsv
\end{cmd}

Columns: \texttt{chrom, start, end, group.a, group.b, estimated.identity,
group.a.length, group.b.length}. The \texttt{group} columns are haplotype
identifiers in \texttt{<sample>\#<hap>} format. About 1.4M pair-window rows.

\subsection*{Run the explorer}

\begin{cmd}
\$ python3 code/01\_explore\_ibs.py
\end{cmd}

This loads the TSV, splits the rows into intra-individual pairs
(e.g.\ \texttt{NA12879\#1} vs.\ \texttt{NA12879\#2}) and inter-individual
pairs (everything else), and writes
\texttt{figures/student\_ibs\_histogram.png}. Open it.

\begin{quest}[1.1]
The right tail of the histogram concentrates at identity $\ge 0.9995$.
What biological scenarios produce values that high? Name two distinct
ones.
\end{quest}

\begin{quest}[1.2]
In this pedigree-only dataset, the intra-individual and inter-individual
means are extremely close, and the inter-individual mean is in fact
\emph{slightly higher}. That is counter-intuitive. Why does it make sense
for this cohort? \emph{Background you need:} the intra-individual class
is hap1-vs-hap2 of the \emph{same} person --- two largely unrelated
homologs. The inter-individual class includes every sibling--sibling and
sibling--grandparent pair in the cohort, and those pairs share large IBD
tracts ($\sim$50\% and $\sim$25\% respectively). Use that to explain the
direction of the gap.
\end{quest}

\begin{quest}[1.3]
The histogram is plotted on \texttt{[0.985, 1.001]}, but the raw
\texttt{estimated.identity} column contains values as low as $\sim$0.12.
From what you know about pangenome alignment to a linear reference,
what produces values that low? (Hint: these are not biology; they are
technical.)
\end{quest}

\section{Part 2 --- IBD detection (12 min)}

\begin{warn}[If ``HMM'' is new]
The IBD detector is a 2-state hidden Markov model: each window is in
state \texttt{IBD} or \texttt{not-IBD}, transitions between adjacent
windows are calibrated by recombination rate, and the emission
$P(\text{identity}\mid\text{state})$ is learned by Baum--Welch from the
data. The flags below (\texttt{--p-enter-ibd},
\texttt{--baum-welch-iters}) directly control these two HMM ingredients.
See the talk slide ``HMM in 30 seconds'' for the picture.
\end{warn}

\subsection*{What we expect to find}

Full siblings share, on average, 50\% of the genome IBD. On a single 90 Mb
arm we therefore expect several multi-Mb IBD segments per sibling pair.
Between \textbf{each sibling and their paternal grandparents}
(NA12889/NA12890) the expected fraction is 25\%; same between each
sibling and their maternal grandparents (NA12891/NA12892). Between the
two grandparental couples, expected IBD is essentially zero.

\subsection*{Run the IBD detector}

\begin{cmd}
\$ bash code/02\_run\_ibd.sh
\end{cmd}

This invokes \texttt{bin/ibd} on the pedigree identity matrix and writes
\texttt{solutions/ibd\_segments.tsv}. The relevant flags:

\begin{cmd}
--similarity-file data/identity\_chr12\_pedigree.tsv \\
--region chr12:40000001-130000000 \\
--size 10000 \\
--min-len-bp 2000000 \\
--min-lod 3.0 \\
--baum-welch-iters 20
\end{cmd}

\begin{itemize}\setlength{\itemsep}{0.2em}
  \item \texttt{--min-len-bp 2000000}: minimum reported segment length
    (2 Mb). Standard value for an arm-scale scan.
  \item \texttt{--min-lod}: minimum LOD score
    $\log_{10}\!\bigl[P(\text{data}\mid\text{IBD}) \,/\, P(\text{data}\mid\neg\text{IBD})\bigr]$.
    A value of 3 means the segment is $10^3\times$ more likely under IBD.
  \item \texttt{--baum-welch-iters 20}: EM iterations for emission
    parameter learning.
\end{itemize}

\subsection*{Read the output}

\begin{cmd}
\$ head -20 solutions/ibd\_segments.tsv
\end{cmd}

Columns: \texttt{chrom, start, end, group.a, group.b, n\_windows,
mean\_identity, mean\_posterior, ..., lod\_score}. Each row is one
detected IBD segment between one haplotype pair.

\begin{quest}[2.1]
Find one IBD segment between a pair of siblings (any of NA12879, NA12881,
NA12882, NA12885, NA12886 with another). Write down the interval and the
LOD score. Bonus: find a pair with more than one multi-Mb IBD segment on
this arm --- how many breakpoints are implied?
\end{quest}

\begin{quest}[2.2]
There are detected ``IBD'' segments between \texttt{NA12892\#1} and
\texttt{NA12892\#2} --- the two haplotypes of one person --- and also
between the two haplotypes of \emph{other} grandparents and even of some
siblings. Strict IBD between the two haplotypes of one diploid individual
would imply autozygosity. What conservative explanation should you
consider first, and what would you check to discriminate?
\end{quest}

\begin{quest}[2.3]
The HMM has two key transition parameters that you did not tune:
\texttt{--p-enter-ibd} (default $10^{-4}$) and
\texttt{--expected-seg-windows} (default 50). What would you expect to
happen to the count of detected segments if you set
\texttt{--p-enter-ibd} to $10^{-2}$? What about
\texttt{--expected-seg-windows} to 5?
\end{quest}

\section{Part 3 --- Local ancestry inference (16 min)}

\subsection*{How the chimeras and the panel were built}

The five chimeric query haplotypes were stitched from real HPRCv2
assemblies. Each \texttt{CHIM\_NN\#1} is a mosaic of one AFR donor and
one EUR donor, copied along ground-truth tracts:

\medskip
\begin{tabular}{@{}llll@{}}
\toprule
chimera & AFR donor & EUR donor & ground-truth tracts \\
\midrule
CHIM\_00\#1 & HG01884 & HG00097 & \texttt{data/ground\_truth\_tracts.tsv} \\
CHIM\_01\#1 & HG01884 & HG00099 & \texttt{...} \\
CHIM\_02\#1 & HG01884 & HG00126 & \texttt{...} \\
CHIM\_03\#1 & HG01884 & HG00128 & \texttt{...} \\
CHIM\_04\#1 & HG01884 & HG00133 & \texttt{...} \\
\bottomrule
\end{tabular}

\medskip

The reference panel in \texttt{data/populations.tsv} is a different
set of 20 HPRCv2 haplotypes: 10 AFR (HG02257, HG02258, HG02280,
HG02451, HG02572, both \texttt{\#1} and \texttt{\#2}) and 10 EUR
(HG00140, HG00146, HG00232, HG00253, HG00272, both haps each).
\textbf{None of the panel samples are donors of the chimeras.} It is
a fair test: the model has to recognize the donor's population
without ever having seen the donor.

\subsection*{The dataset}

\texttt{data/identity\_chr12\_admix.tsv} contains pairwise identity
per 10-kb window for \texttt{chr12:50--120\,Mb} --- 70 Mb in the
middle of the long q arm. Because 70 Mb is much longer than a
typical admixture tract, each chimera carries 3--7 ground-truth
breakpoints in this window.

\subsection*{Run local ancestry inference}

\begin{cmd}
\$ bash code/03\_run\_ancestry.sh
\end{cmd}

The relevant flags:

\begin{cmd}
--similarity-file data/identity\_chr12\_admix.tsv \\
--query-samples data/queries.txt \\
--populations data/populations.tsv \\
--region chr12:50000001-120000000 \\
--window-size 10000 \\
--auto-configure \\
--identity-floor 0.9
\end{cmd}

\begin{itemize}\setlength{\itemsep}{0.2em}
  \item \texttt{--auto-configure}: estimates two emission-model
    knobs from the data, instead of asking you to set them.
    \begin{itemize}\setlength{\itemsep}{0.1em}
      \item \emph{Softmax temperature $T$}: rescales per-window
        similarity differences before turning them into emission
        probabilities via $\text{softmax}(\Delta/T)$. Small $T$ $\to$
        sharp commit to one population from tiny similarity gaps;
        large $T$ $\to$ smooth, conservative posteriors.
      \item \emph{Pairwise-contrast weight}: how much weight the
        emission gives to the \emph{difference} between
        max-similarity-to-population-A and max-similarity-to-B,
        relative to the absolute identity. High contrast weight
        emphasizes which panel is closer; low contrast weight makes
        the emission depend more on raw identity values.
    \end{itemize}
  \item \texttt{--identity-floor 0.9}: drops windows whose maximum
    similarity across the panel is below 0.9 --- usually alignment-gap
    artifacts.
\end{itemize}

\subsection*{Inspect the painting}

\begin{cmd}
\$ head -20 solutions/ancestry\_segments.tsv \\
\$ python3 code/04\_plot\_painting.py
\end{cmd}

This produces \texttt{figures/student\_ancestry\_painting.png} with two
bars per chimera: ``truth'' on top, ``impop'' below. Open it.

\begin{quest}[3.1]
Pick one chimera and count its ground-truth breakpoints in the 50--120 Mb
interval. How many of them does the model recover within $\sim$1 Mb?
Which ones does it miss or hallucinate?
\end{quest}

\begin{quest}[3.2]
Scan the five paintings side by side. Are the errors \emph{the same kind}
across chimeras, or chimera-specific? A recurring error in roughly the
same genomic region across multiple, independent chimeras suggests a
property of the \emph{region}, not of the chimeras. Do you see one?
\end{quest}

\begin{quest}[3.3]
Look at the \texttt{mean\_posterior} column of the TSV. Several wrong
calls carry posteriors $\ge 0.9$. How can the model be confidently wrong?
\emph{Hint: think about the difference between the model being
confident in its assumptions and the model being correct against reality.}
\end{quest}

\subsection*{Open the black box: forward-backward, Viterbi, and the raw input}

The TSV you just read is the \emph{Viterbi} decoding: the single most
likely sequence of ancestry states. The HMM also computes a
\emph{marginal posterior} per window via forward-backward: the
probability of each state at each window, integrating over all paths.
The script invoked \texttt{ancestry} with \texttt{--posteriors-output},
so you already have these per-window numbers in
\texttt{solutions/ancestry\_posteriors.tsv}.

\begin{cmd}
\$ head -3 solutions/ancestry\_posteriors.tsv \\
\$ python3 code/07\_plot\_posteriors.py
\end{cmd}

This produces two figures:

\begin{itemize}\setlength{\itemsep}{0.2em}
  \item \texttt{figures/student\_posterior\_trace.png} for
    \texttt{CHIM\_03\#1}: ground truth on top, P(AFR) per window in
    the middle (the marginal posterior), Viterbi-decoded segments at
    the bottom. The dashed line at $P=0.5$ marks where Viterbi flips.
  \item \texttt{figures/student\_panel\_similarity.png}: the raw
    input the HMM emissions see. Per window, $\Delta$ = (max identity
    to AFR panel) $-$ (max identity to EUR panel), smoothed over
    250\,kb. Notice how small the contrast is --- median $|\Delta|
    \approx 1$--$3 \times 10^{-4}$ per single window. The HMM has to
    chain dozens of windows to make a confident call from such a
    weak per-window signal.
\end{itemize}

\begin{quest}[3.4]
Open \texttt{student\_posterior\_trace.png}. Find a window where
P(AFR) sits between 0.4 and 0.6: the model is genuinely uncertain
there. What is Viterbi's call at that same window, and what does
that tell you about the difference between the marginal posterior
(forward-backward) and the most-likely-path (Viterbi)?
\end{quest}

\section{Part 4 --- Grandparent painting of the CEPH pedigree (7 min)}

We swap roles. Instead of a population panel, we use the \emph{four
grandparents} as four ``populations'' (K=4), and paint each
grandchild's two haplotypes against them. This is the same HMM, same
binary, same identity matrix as Part 2 --- only the panel file changes.

\subsection*{The setup}

\begin{cmd}
\$ cat data/pedigree\_populations.tsv \\
\$ cat data/pedigree\_queries.txt
\end{cmd}

Four ``populations'', each containing the two haplotypes of one
grandparent:

\begin{cmd}
GP\_PAT  NA12889\#1, NA12889\#2  (paternal grandfather) \\
GM\_PAT  NA12890\#1, NA12890\#2  (paternal grandmother) \\
GP\_MAT  NA12891\#1, NA12891\#2  (maternal grandfather) \\
GM\_MAT  NA12892\#1, NA12892\#2  (maternal grandmother)
\end{cmd}

The ten queries are the ten haplotypes of the five grandchildren
(NA12879, NA12881, NA12882, NA12885, NA12886 --- \texttt{\#1} and
\texttt{\#2} each).

\subsection*{Run the painting}

\begin{cmd}
\$ bash code/05\_paint\_pedigree.sh \\
\$ python3 code/06\_plot\_pedigree\_painting.py
\end{cmd}

This produces \texttt{figures/student\_pedigree\_painting.png}: two
stacked bars per grandchild, one per haplotype, coloured by inferred
grandparent of origin.

\subsection*{What you should see}

The opening slide of the talk (``a chromosome inherited through a
pedigree'') is now a concrete result on real data. In particular:

\begin{itemize}\setlength{\itemsep}{0.2em}
  \item Each grandchild's two haplotypes split cleanly: one is a mosaic
    of \emph{only} GP\_PAT + GM\_PAT blocks (paternal homolog), the
    other of \emph{only} GP\_MAT + GM\_MAT blocks (maternal homolog).
  \item Breakpoints within a homolog mark actual meiotic recombinations
    in the mother or father of the CEPH family.
  \item The VCF never told us which haplotype index (\texttt{\#1} or
    \texttt{\#2}) was paternal. The painting does.
\end{itemize}

\begin{quest}[4.1]
For one grandchild of your choice, tell me which of
\texttt{\#1} and \texttt{\#2} is paternal and which is maternal. Justify
from the figure, not from the FASTA headers.
\end{quest}

\begin{quest}[4.2]
Pick a grandchild with $\ge 2$ breakpoints on one homolog in the
40--130 Mb window. How many generations of meiotic recombination do
those breakpoints represent (think about which parent generated them)?
\end{quest}

\begin{quest}[4.3]
If the model ever placed a GP\_MAT block on the same homolog as a
GP\_PAT block, that would be biologically impossible --- a grandchild
cannot receive one homolog from both father and mother. Do you see any
such violations in your figure? If yes, what would be the first thing
you would check before blaming the biology?
\end{quest}

\section*{Hand-in (optional)}

If your instructor asks: write a 1-paragraph summary covering
(a) one thing you found in the IBS histogram that was not obvious before,
(b) one IBD segment you can confidently call between two named
haplotypes, (c) one ancestry call from Part 3 where the model is wrong
with high confidence and what that tells you about calibration, and
(d) which grandchild haplotype indices (\texttt{\#1} vs \texttt{\#2})
you concluded to be paternal vs maternal from Part 4.

\end{document}
